% Using InitechBase Chapter 5, entries ? .. DO WHILE (AN-04).

% ? -- os/samir/cmd/query.c q_cmd_hook "?" (leading CR LF, values joined by
% one blank, no trailing newline; q_render_val: N without trailing zeros of
% a computed value, field N without its decimals, L as .T./.F., D by SET
% DATE).  Gates: test-samir-query, test-interp-list.
\begin{XbaseEntry}{?}{command}
Displays the value of one or more expressions on a new line.
\Syntax
\SyntaxLine{\kw{?} \opt{\xp{expression list}}}
\Usage
\cmd{?} moves to the beginning of a new line and displays the value of each
expression in the list.  The values are separated by one blank.  With no
expression, \cmd{?} displays an empty line.

Character values are displayed with their trailing blanks, so a character
field takes its full width.  Numbers are displayed without leading blanks;
a numeric field or a whole number is displayed without its decimal places,
and the result of a division or other calculation with as many decimal places
as SET DECIMALS gives (two when InitechBase starts).  Logical values are
displayed as \texttt{.T.} or \texttt{.F.}; dates in the form set by SET DATE.

The cursor stays at the end of the displayed line.  At the dot prompt the next
prompt therefore follows on the same line as the value, as the examples show.
\Tips
Use \cmd{?} at the dot prompt to look at a field of the current record, a
memory variable or the result of a function without listing the record.
\Examples
\DotExample{examples/c-qmark}
A character string and a number:
\DotExample{examples/c-qmark-2}
\SeeAlso{??, DISPLAY, LIST, SET DECIMALS, STR(), TRANSFORM()}
\end{XbaseEntry}

% ?? -- os/samir/cmd/query.c "??" (no leading CR LF).  Gate: test-samir-query.
\begin{XbaseEntry}{??}{command}
Displays the value of one or more expressions where the cursor is, without
moving to a new line.
\Syntax
\SyntaxLine{\kw{??} \xp{expression list}}
\Usage
\cmd{??} displays the values in the same way as \cmd{?}, but at the current
position of the cursor: it does not begin a new line.  A program uses it to
build one line of output from several commands.
\Examples
\DotExample{examples/c-qqmark}
\SeeAlso{?, DISPLAY, LIST}
\end{XbaseEntry}

% @ SAY -- say-get-read.md Sections 1-2; programming-and-io.md Section 7;
% real III+ fns-mint/{row,col,prow,pcol}/GOLD.OUT.
\begin{XbaseEntry}{@}{command}
Displays an expression at a specified row and column of the current device.
\Syntax
\SyntaxLine{\kw{@} \xp{row}\kw{,}\xp{column} \kw{SAY} \xp{expression}}
\Usage
Rows and columns are numbered from zero.  On the screen, rows run from 0
through 24 and columns from 0 through 79.  The coordinates are numeric
expressions; a fractional part is dropped.  A screen position outside those
limits raises \xmsg{30}{Position is off the screen.}.
This version accepts the same coordinate limits for printer output; a
printer position beyond those limits also raises error 30.

The expression is displayed with the same formatting as \cmd{?}.  The cursor
then stands immediately after the displayed value.  SET DEVICE TO PRINT
sends the positioned output to the printer instead, moving its print head.
Positioned output is not copied to the SET ALTERNATE file.

This version accepts the SAY form shown above.  A PICTURE clause raises
\xmsg{36}{Unrecognized phrase/keyword in command.}.
\Examples
Save the cursor position before displaying it with \cmd{?}:
\DotExample{examples/cmd-say}
\SeeAlso{SET DEVICE, ROW(), COL(), PROW(), PCOL(), ?}
\end{XbaseEntry}

% ACCEPT -- os/samir/cmd/proc.c p_accept (prompt expression, then one
% cooked line into a Character memvar; empty line -> "").  Gate:
% test-interp-proc.
\begin{XbaseEntry}{ACCEPT}{command}
Displays a prompt, reads a line typed at the keyboard and stores it in a
memory variable as a character string.
\Syntax
\SyntaxLine{\kw{ACCEPT} \opt{\xp{expC}} \kw{TO} \xp{memvar}}
\Usage
\cmd{ACCEPT} displays the character expression \xp{expC}, if one is given, and
waits for a line to be typed.  The line ends when you press \key{Enter}.  The
characters typed are stored in \xp{memvar} as a character string, whatever
they are: a number typed in answer to \cmd{ACCEPT} is stored as characters.  If
you press \key{Enter} alone, \xp{memvar} is an empty string.  The memory
variable is created if it does not exist.

The prompt is usually a character string in quotation marks.  It is displayed
as it stands, so end it with a blank or a colon and a blank to keep the answer
apart from it.
\Tips
Use \cmd{INPUT} when the answer is a number, a date or a logical value, and
\cmd{WAIT} when one keystroke is enough.  To compare an answer with stored
data in capitals, convert it with \fn{UPPER()}.
\Examples
\DotExample{examples/c-accept}
Without a prompt, \cmd{ACCEPT} waits on a new line:
\DotExample{examples/c-accept-2}
\SeeAlso{INPUT, STORE, UPPER(), WAIT}
\end{XbaseEntry}

% APPEND BLANK -- os/samir/cmd/mutate.c m_append (BLANK required; adds a
% blank record, makes it current, inserts its key into every open index).
% Gate: test-interp-replace, test-samir-write.
% APPEND FROM -- os/samir/cmd/mutate.c d_append_from (Using III+ U5-28/PDF
% 140: fields matched by name in any order; deleted records skipped only
% under SET DELETED ON; default .DBF; the database in USE refused, 3; a
% memo field in the database in USE, 34; a field of the same name and
% another type, 9; every open index rebuilt).  Lane samirdml 14a8fc75.
% Gates: test-samir-dml, test-samir-dml-emu.
\begin{XbaseEntry}{APPEND}{command}
Adds records to the end of the database in USE: one blank record, or the
records of another database.
\Syntax
\SyntaxLine{\kw{APPEND BLANK}}
\par or\par
\SyntaxLine{\kw{APPEND FROM} \xp{file}}
\Defaults
\kw{APPEND FROM} assumes the extension \file{.DBF}.
\Usage
\cmd{APPEND BLANK} adds one record to the end of the database in the selected
work area and makes it the current record.  Every field of the new record is
empty: character fields are blanks, numeric fields are zero, date fields are
blank and logical fields are false.  The record count, \fn{RECCOUNT()}, grows
by one.

The new record is written to the disk at once.  If index files are open with
the database, the new record is entered in each of them.

Fill in the fields of the new record with \cmd{REPLACE}.  In a program, an
\cmd{ACCEPT} or \cmd{INPUT} for each field followed by \cmd{APPEND BLANK} and
\cmd{REPLACE} enters a record.

\cmd{APPEND FROM} adds every record of the database named to the end of the
database in USE, in the order of its records.  Fields are matched by name:
a field of the database in USE takes the value of the field of the same name
in the other database, wherever it stands in its structure.  A field the
other database does not have is left empty, and a field the database in USE
does not have is not copied.  A character value longer than its field is cut
off on the right.  Records marked for deletion are appended, and are not
marked in the database in USE, unless SET DELETED is ON: then they are left
out.  The last record appended becomes the current record, and every open
index file is brought up to date.

The database named must not be the database in USE
(\xmsg{3}{File is already open.}), and two fields of the same name must have
the same type (\xmsg{9}{Data type mismatch.}).  Records cannot be appended to
a database that has a memo field (\xmsg{34}{Operation with Memo field
invalid.}).
\Examples
A new record is added and two of its fields are filled in:
\DotExample{examples/c-append}
The batch of new invoices on the sample diskette, \file{ARBATCH.DBF}, has the
fields of the ledger in another order, and no \texttt{PAID} field.  Its three
records are added to \file{INVOICE.DBF}:
\DotExample{examples/c-append-2}
\SeeAlso{REPLACE, RECCOUNT(), DELETE, COPY FILE}
\end{XbaseEntry}

% CANCEL -- os/samir/cmd/proc.c "CANCEL" (INTERP_RC_CANCEL unwinds every DO
% level to the dot prompt; Using III+ U5-43).  Gate: test-samir-ctl.
\begin{XbaseEntry}{CANCEL}{command}
Ends the program and returns to the dot prompt.
\Syntax
\SyntaxLine{\kw{CANCEL}}
\Usage
In a program, \cmd{CANCEL} stops the program at once, together with every
program and procedure that called it, and the dot prompt returns.  Memory
variables made \cmd{PUBLIC} remain.  Typed at the dot prompt, \cmd{CANCEL}
does nothing, except while a program is suspended (see \cmd{RESUME}): then it
ends the suspended program.
\SeeAlso{QUIT, RESUME, RETURN, SET ESCAPE}
\end{XbaseEntry}

% CLEAR MEMORY -- os/samir/cmd/flow.c exec_command "CLEAR MEMORY"
% (xb_interp_release_all: every memory variable, PUBLIC and PRIVATE, at any
% level; Using III+ U5-51/PDF 163).  Lane samirdml 14a8fc75.  Gates:
% test-samir-dml, test-samir-dml-emu.
\begin{XbaseEntry}{CLEAR MEMORY}{compound}\XAlso{CLEAR}{CLEAR}%
Releases all memory variables.
\Syntax
\SyntaxLine{\kw{CLEAR MEMORY}}
\Usage
\cmd{CLEAR MEMORY} releases every memory variable: the public ones, and the
private ones of every program and procedure running.  A memory variable
released no longer exists; a reference to it gives
\xmsg{12}{Variable not found.}  Databases and other files are not affected.

At the dot prompt \cmd{CLEAR MEMORY} does what \cmd{RELEASE ALL} does.  In a
program they differ: \cmd{RELEASE ALL} releases only the private memory
variables of the routine running, \cmd{CLEAR MEMORY} all of them.
\Examples
\DotExample{examples/c-clearmem}
\SeeAlso{RELEASE, PUBLIC, PRIVATE, SAVE}
\end{XbaseEntry}

% CLOSE -- os/samir/samir_main.c repl_do_close (Using III+ U5-53/PDF 165:
% CLOSE <file type> / ALL.  ALL and DATABASES close every database with its
% indexes; INDEX closes the index files of every work area and leaves the
% databases open in natural order; ALTERNATE, FORMAT and PROCEDURE are
% accepted and close nothing -- SET ALTERNATE TO closes the alternate file,
% observed in the replay; any other word, or none, is 10).  Lane samirdml 14a8fc75.
% Gates: test-samir-repl, test-samir-dml, test-samir-dml-emu.
\begin{XbaseEntry}{CLOSE}{command}
Closes the open files of one type, or all of them.
\Syntax
\SyntaxLine{\kw{CLOSE} \xp{file type}}
\par or\par
\SyntaxLine{\kw{CLOSE ALL}}
\Usage
\cmd{CLOSE} is followed by one of the words \kw{ALL}, \kw{DATABASES},
\kw{INDEX}, \kw{ALTERNATE}, \kw{FORMAT} or \kw{PROCEDURE}.  Memory variables
are not affected by any of them.

\cmd{CLOSE ALL} and \cmd{CLOSE DATABASES} close the database open in every
work area, with its index files.  Afterwards no database is in USE in any
work area.

\cmd{CLOSE INDEX} closes the index files of every work area and leaves the
databases open.  Each database is then read in the order of its records
until an index file is opened again with \cmd{SET INDEX} or \cmd{USE}.

\cmd{CLOSE ALTERNATE}, \cmd{CLOSE FORMAT} and \cmd{CLOSE PROCEDURE} are
accepted and close nothing in this version.  Close the alternate file with
\cmd{SET ALTERNATE TO}, without a file name.  \cmd{CLOSE} without a word, or with any other word,
is reported as \xmsg{10}{Syntax error.}
\Tips
To close only the database of the selected work area, type \cmd{USE} with no
file name.
\Examples
\DotExample{examples/c-close}
With its index file closed, the database is read in the order of its
records again:
\DotExample{examples/c-close-2}
\SeeAlso{QUIT, SELECT, USE, SET INDEX}
\end{XbaseEntry}

% CONTINUE -- os/samir/cmd/query.c q_continue (re-applies the stored LOCATE
% FOR from the record after the current one; #42 without LOCATE).
% Gate: test-samir-query.
\begin{XbaseEntry}{CONTINUE}{command}
Moves to the next record that meets the condition of the last \cmd{LOCATE}.
\Syntax
\SyntaxLine{\kw{CONTINUE}}
\Usage
\cmd{CONTINUE} goes on with the search begun by \cmd{LOCATE} in the selected
work area.  It tests the records after the current one against the
\kw{FOR} condition of the \cmd{LOCATE} and stops at the first that meets it.
\fn{FOUND()} is then true.  When no further record meets the condition, the
record pointer is at the end of the file: \fn{EOF()} is true and
\fn{FOUND()} false.

\cmd{CONTINUE} uses only the \kw{FOR} condition.  The scope and the \kw{WHILE}
condition of the \cmd{LOCATE}, if it had them, no longer apply.  Without an
earlier \cmd{LOCATE}, \cmd{CONTINUE} displays
\xmsg{42}{CONTINUE without LOCATE.}
\Tips
In a program, test \fn{FOUND()} after each \cmd{CONTINUE}, and stop when it is
false.
\Examples
\DotExample{examples/c-continue}
\SeeAlso{LOCATE, FOUND(), EOF(), SEEK}
\end{XbaseEntry}

% COPY FILE -- os/samir/cmd/mutate.c d_file (Using III+ U5-59/PDF 171: any
% closed file, both names with their extensions; a source open as a
% database, 3; absent, 1; a target open as a database, 3; SAFETY ON
% refuses an existing target, 7; real DML-11: the copy is written in
% 512-byte blocks, the last filled out with zero bytes).  Lane samirdml
% 14a8fc75.  Gates: test-samir-dml, test-samir-dml-emu.
\begin{XbaseEntry}{COPY FILE}{compound}\XAlso{COPY}{COPY}%
Makes a copy of a file.
\Syntax
\SyntaxLine{\kw{COPY FILE} \xp{file} \kw{TO} \xp{new file}}
\Usage
\cmd{COPY FILE} copies the first file to the second.  Give both names with
their extensions; none is assumed.  Any file may be copied: a database, an
index file, a program, a report form.  The copy is written in blocks of 512
bytes, and the last block is filled out with zero bytes.

A database that is open cannot be copied, and cannot be replaced by the copy:
\cmd{COPY FILE} reports \xmsg{3}{File is already open.}  If the first file is
not on the diskette, it reports \xmsg{1}{File does not exist.}  While SET
SAFETY is ON, a file of the new name is not replaced
(\xmsg{7}{File already exists.}).

A database's memo fields are kept in its memo file (\file{.DBT}).  Copy the
memo file with the database, under the same new name.
\Examples
A copy of the departments database is made and opened:
\DotExample{examples/c-copyfile}
\SeeAlso{ERASE, RENAME, TYPE}
\end{XbaseEntry}

% DELETE -- os/samir/cmd/mutate.c m_delete (default the current record;
% FOR without a scope = ALL; marks with 0x2A, written at once).
% Gate: test-interp-replace, test-dbf-mutate.
\begin{XbaseEntry}{DELETE}{command}
Marks records of the database in USE for deletion.
\Syntax
\SyntaxLine{\kw{DELETE} \opt{\xp{scope}} \opt{\kw{FOR} \xp{condition}}
  \opt{\kw{WHILE} \xp{condition}}}
\Defaults
Without a scope or a condition, \cmd{DELETE} marks the current record only.
With a \kw{FOR} condition and no scope, the scope is \kw{ALL}.
\Usage
A record marked for deletion stays in the database.  \cmd{LIST} and
\cmd{DISPLAY} show it with an asterisk in front of its record number, and the
function \fn{DELETED()} is true for it.  \cmd{RECALL} removes the mark;
\cmd{PACK} removes the marked records from the database for good.

While SET DELETED is ON, the commands that process records leave out the
records marked for deletion.

The mark is written to the disk at once.
\Examples
\DotExample{examples/c-delete}
\SeeAlso{PACK, RECALL, SET DELETED, ZAP, DELETED()}
\end{XbaseEntry}

% DISPLAY -- os/samir/cmd/query.c q_list (is_display: default scope the
% current record; heading line per SET HEADING; deleted records '*').
% Gates: test-samir-query, test-samir-list-layout, test-interp-list.
\begin{XbaseEntry}{DISPLAY}{command}
Displays the current record, or the records in a scope.
\Syntax
\SyntaxLine{\kw{DISPLAY} \opt{\xp{scope}} \opt{\xp{expression list}}
  \opt{\kw{FOR} \xp{condition}} \opt{\kw{WHILE} \xp{condition}}}
\Defaults
The scope is the current record.  Without an expression list, every field of
the record is displayed.
\Usage
\cmd{DISPLAY} works as \cmd{LIST} does, with one difference: without a scope it
shows only the current record.  With a scope, such as \kw{ALL} or \kw{NEXT}
\xp{n}, it shows the records of the scope that meet the conditions.  A
\kw{FOR} condition alone tests the current record only; to search the
database, give the scope \kw{ALL}.

The first line is a heading (unless SET HEADING is OFF): \texttt{Record\#}
and the names of the fields, or the expressions, as typed.  Each record follows
on its own line, with its record number.  A record marked for deletion has an
asterisk in front of its number.  A memo field is displayed with its text.

The scope must come before the expression list; \kw{FOR} and \kw{WHILE} may
follow it.  A \cmd{DISPLAY} of a scope leaves the record pointer on the last
record it examined.
\Examples
\DotExample{examples/c-display}
Every record that meets a condition:
\DotExample{examples/c-display-2}
\SeeAlso{LIST, ?, SET HEADING, GO, SKIP}
\end{XbaseEntry}

% Ground truth: Using III+ U5-108/109; real inspection mint.
% Gates: test-samir-dml-memory/status, their -emu and mutant legs.
% A form of the DISPLAY command (compound): its NAME is registered by the
% DISPLAY entry; the inventory keys verbs, as for CLEAR MEMORY (initech-u01i).
\begin{XbaseEntry}{DISPLAY MEMORY}{compound}
Shows the memory variables and pauses for a key between screens.
\Syntax
\SyntaxLine{\kw{DISPLAY MEMORY}}
\Usage
Each entry gives the name, PUBLIC or PRIVATE status, type, and current value.
Hidden outer variables are included and marked hidden.  A PRIVATE declaration
has no entry until a value is assigned; a newly declared PUBLIC is logical
false.  Numbers appear both at their display precision and with eight decimal
places.  PRIVATE entries name the program that created them.

The totals give the variables and bytes used and available.  There are 256
variable slots and 6,000 bytes for variable values.  No database need be open.
Inspecting memory preserves the variables and the selected work area.
\SeeAlso{LIST MEMORY, PUBLIC, PRIVATE, STORE, RELEASE}
\end{XbaseEntry}

% A form of the DISPLAY command (compound): its NAME is registered by the
% DISPLAY entry; the inventory keys verbs, as for CLEAR MEMORY (initech-u01i).
\begin{XbaseEntry}{DISPLAY STATUS}{compound}
Shows the open files and working environment, pausing for a key between pages.
\Syntax
\SyntaxLine{\kw{DISPLAY STATUS}}
\Usage
Each open work area shows its database name and alias, memo file, indexes and
their keys, controlling index, filter, and configured relation.  The selected
area is identified.  The report also shows the alternate file, default drive,
ON ERROR and ON ESCAPE commands, SET states, and function key assignments.
It may be used with no database open and does not move any record pointer or
change the selected work area.  The relation shown is its configured setting;
automatic movement of a related area's record pointer is not provided.
\SeeAlso{LIST STATUS, SELECT, SET INDEX, SET FILTER, SET RELATION}
\end{XbaseEntry}

% DO -- os/samir/samir_main.c repl_try_do_file (DO <name> at the dot prompt
% runs <name>.PRG from the disk; 8191-byte cap, #18) and os/samir/cmd/proc.c
% p_do / p_call (a PROCEDURE of the program running; WITH arguments by value;
% #16 unknown target).  Gates: test-samir-canon-y2k, test-interp-proc.
\begin{XbaseEntry}{DO}{command}
Runs a program file, or a procedure of the program that is running.
\Syntax
\SyntaxLine{\kw{DO} \xp{program}}
\par or\par
\SyntaxLine{\kw{DO} \xp{procedure} \opt{\kw{WITH} \xp{expression list}}}
\Usage
At the dot prompt, \cmd{DO} \xp{program} reads the program file
\xp{program}\texttt{.PRG} from the disk and runs it.  You need not type the
extension.  The program runs from its first line until it reaches
\cmd{RETURN}, its first \cmd{PROCEDURE} or its last line; then the dot prompt
returns.  A program file may hold at most 8,191 characters; a larger file is
refused with \xmsg{18}{Line exceeds maximum of 254 characters.}

In a program, \cmd{DO} \xp{procedure} runs a procedure of the same program
file: the lines that follow its \cmd{PROCEDURE} line, up to the next
\cmd{RETURN}.  When the procedure returns, the program goes on with the line
after the \cmd{DO}.  Procedures may call other procedures, and themselves.  At
most 20 \cmd{DO} commands may be in progress at once, counting the program
started at the dot prompt; the \cmd{DO} that would begin a twenty-first is
refused with \xmsg{103}{DOs nested too deep.}

\kw{WITH} passes the values of the expressions to the procedure, which
receives them in the memory variables of its \cmd{PARAMETERS} line, in order.
The procedure receives a copy of each value: a change it makes to a parameter
does not change the expression that was passed.
\draft{a character constant written directly in the WITH list is not passed
intact by this version (proc.c); pass a character value
through a memory variable or an expression}

A name that is neither a program file on the disk nor a procedure of the
program running is refused with \xmsg{16}{*** Unrecognized command verb.}
\Tips
The programs that come with InitechBase work on the database already in USE:
open it at the dot prompt with \cmd{USE} before you type \cmd{DO}.
\Examples
\file{RAISE.PRG} calls its procedure \texttt{NEWSAL} with two values:
\ProgramListing{RAISE.PRG}
\DotExample{examples/c-do}
\SeeAlso{PARAMETERS, PROCEDURE, RETURN, PUBLIC, PRIVATE}
\end{XbaseEntry}

% DO CASE -- os/samir/cmd/flow.c LK_DOCASE/LK_CASE/LK_OTHERWISE/LK_ENDCASE
% (first true CASE runs; guard must be Logical, #37; program lines only).
% Gate: test-interp-flow.
\begin{XbaseEntry}{DO CASE}{compound}
\XAlso{CASE}{CASE}\XAlso{OTHERWISE}{OTHERWISE}\XAlso{ENDCASE}{ENDCASE}%
Chooses one of several groups of commands in a program.
\Syntax
\SyntaxLine{\kw{DO CASE}}
\SyntaxLine{\kw{CASE} \xp{condition}}
\SyntaxLine{\ \ \xp{commands}}
\SyntaxLine{\opt{\kw{CASE} \xp{condition}}}
\SyntaxLine{\ \ \opt{\xp{commands}}}
\SyntaxLine{\rep}
\SyntaxLine{\opt{\kw{OTHERWISE}}}
\SyntaxLine{\ \ \opt{\xp{commands}}}
\SyntaxLine{\kw{ENDCASE}}
\Usage
InitechBase tests the condition of each \cmd{CASE} in turn.  It carries out
the commands of the first \cmd{CASE} whose condition is true, and then goes on
after \cmd{ENDCASE}; the other \cmd{CASE} groups are skipped, even if their
conditions are also true.  If no condition is true, the commands after
\cmd{OTHERWISE} are carried out; without \cmd{OTHERWISE}, none are.

Each condition must be a logical expression.  A condition of another type
stops the program with \xmsg{37}{Not a Logical expression.}

\cmd{DO CASE} is used only in programs.  The commands in each group may be any
commands, including further \cmd{DO CASE}, \cmd{IF} and \cmd{DO WHILE}
structures.
\Tips
Use \cmd{DO CASE} instead of a chain of \cmd{IF} commands whenever there are
more than two possibilities.
\Examples
The program \file{DEPTNAME.PRG} gives the name of the department of the
current employee:
\ProgramListing{DEPTNAME.PRG}
\DotExample{examples/c-docase}
\SeeAlso{IF, DO WHILE}
\end{XbaseEntry}

% DO WHILE -- os/samir/cmd/flow.c LK_DOWHILE/LK_ENDDO/LK_LOOP/LK_EXIT (guard
% must be Logical, #37; LOOP back to the test; EXIT after ENDDO).
% Gates: test-interp-flow, test-samir-canon-y2k.
\begin{XbaseEntry}{DO WHILE}{compound}
\XAlso{ENDDO}{ENDDO}\XAlso{LOOP}{LOOP}\XAlso{EXIT}{EXIT}%
Repeats a group of commands in a program as long as a condition is true.
\Syntax
\SyntaxLine{\kw{DO WHILE} \xp{condition}}
\SyntaxLine{\ \ \xp{commands}}
\SyntaxLine{\ \ \opt{\kw{LOOP}}}
\SyntaxLine{\ \ \opt{\kw{EXIT}}}
\SyntaxLine{\kw{ENDDO}}
\Usage
InitechBase tests the condition.  If it is true, the commands up to
\cmd{ENDDO} are carried out and the condition is tested again; when it is
false, the program goes on after \cmd{ENDDO}.  The condition must be a logical
expression, or the program stops with \xmsg{37}{Not a Logical expression.}

\cmd{LOOP} goes back to the test at once, leaving out the rest of the commands
of the loop.  \cmd{EXIT} leaves the loop at once and goes on after
\cmd{ENDDO}.  Both are normally written inside an \cmd{IF}.

The commonest loop of an InitechBase program processes every record of a
database: it begins at the first record, ends when \fn{EOF()} is true, and
moves to the next record with \cmd{SKIP} before \cmd{ENDDO}.

\cmd{DO WHILE} is used only in programs.  At the dot prompt, \cmd{EXIT} is a
different command: it leaves InitechBase (see \cmd{QUIT}).
\Tips
Make sure that something in the loop changes the condition.  A loop that
processes records and forgets the \cmd{SKIP} never ends.
\Examples
\file{PAYTOTAL.PRG} totals the salaries of exempt and non-exempt staff,
leaving out the records marked for deletion with \cmd{LOOP}:
\ProgramListing{PAYTOTAL.PRG}
\DotExample{examples/c-dowhile}
\SeeAlso{DO CASE, IF, SKIP, EOF()}
\end{XbaseEntry}
